\documentclass{article}
\usepackage[utf8]{inputenc}
\usepackage[T1]{fontenc}
\usepackage[margin=2cm]{geometry}
\usepackage{hyperref}
\sloppy
\begin{document}
\section*{Historical source}
Encyclopaedia Britannica scan, source row 756. Parse 3.5 extracted the text. Super 3.5 VL corrected it from an enlarged crop of each block. It can still contain recognition errors.
Page headers, footers, and page numbers are removed.

\subsection*{MAGNETISM (E0069)}
\textit{Continues an entry from an earlier page.}

43
tempered
sheet
from a
the magma
film.

But any degree of temper that is given to steel makes
a very important change in this respect. In the first
place, it acquires magnetism more slowly by induction
than an equal and similar piece of common iron, and
finally acquires less. These differences are easily ex-
plained by the deviations which it causes in the mariner's
needle from the magnetic meridian, and by its attrac-
tion.

When the inducing magnet is removed, some magnetism remains in the fixed bar, which retains the polarity which it had in the neighbourhood of the magnet.

Steel tempered to the degree fit for watch springs acquires a strong magnetism, which it exhibits immediately on the removal of the magnet. But it diffuses very fast; and, in a very few minutes, it is reduced to less than one half of its intensity while in contact with the magnet, and not two thirds of what it was immediately on removal from it. It continues to diffuse for some days, though the bar be kept with care; but the diffusion diminishes fast, and it retains at least one third of its greatest power for any length of time, unless carelessly kept or injudiciously treated.

Steel tempered for strong cutting tools, such as chisels, punches, and drills for metal, acquires magnetism full more slowly by induction, and acquires less of it while in contact with the magnet; but it retains it more firmly, and finally retains a greater proportion of what it had acquired.

Steel made as hard as possible, is much longer in acquiring all the magnetism which simple juxtaposition can give to it. It acquires less than the former; but it retains it with great firmness, and finally retains a much greater proportion.

Such ores of iron as are susceptible of magnetism, are nearly like hard steel in their effects; that is, in the time necessary for their great impregnation, and in the durability of the acquired magnetism. They differ exceedingly in respect to the degree of power which they can attain by mere juxtaposition, and the varieties seem to depend on heterogeneous mixture. We must observe, that few ores of iron are susceptible of magnetism in their natural state. The ordinary ores, consisting of the metal in the state of an oxyd, and combined with sulphur, are not magnetizable while remaining in that state. Most ores require roasting, and a sort of cementation, in contact with inflammable substances. This matter is not well understood; but it would seem that complete metallization is far from being the most favourable condition, and that a certain degree of oxidation, and perhaps some other composition, yet unknown, make the best loadstones. But all this is extremely obscure. The late Dr Gowin Knight made a composition which acquired a very strong and permanent magnetism, but the secret died with him. Dr Gilbert speaks of similar compositions, in which ferruginous clays were ingredients; but we know nothing of the state of the metal in them, nor their mode of acquiring magnetism.

It is of peculiar importance to remark that the acquisition of magnetism is gradual and progressive, and that the gradation is the more perceptible in proportion as the field is of a harder temper. When a magnet is brought to one end of a bar of common iron, its remote extremity, unless exceedingly long, acquires its utmost magnetism immediately. But when the north pole of

a magnet is applied to one end of a bar of hard steel,
the part in contact immediately becomes a fourth pole,
and the far end is not yet affected. We observe a north
pole formed at some distance from the contact, and be-
yond this a faint fourth pole. These gradually advance
along the bar. The remote extremity becomes first a
faint fourth pole, and it is not till after a very long while
(if ever) that it becomes a simple, vigorous, north
pole. More frequently it remains a diffused and feeble
north pole: nay, if the bar be very long, it often hap-
pens that we have a succession of north and fourth poles,
which never make their way to the far end of the bar.
This phenomenon was first observed (we think) by Dr
Brook Taylor, who gives an account of his observations
in the Philosophical Transactions, no 344.

From the account we have given of these phenomena of induced magnetism, it appears that the temporary magnetism is always so disposed that the sum of the mutual attractions of the dissimilar poles exceeds the sum of the repulsions between the similar poles, and that therefore the two magnets tend to each other. This is evidently equivalent to saying, that a piece of unmagnetic iron is always attracted by a magnet. No exception has ever been observed to this fact; for Pliny's story of a Theamedes, or loadstone, which repels iron, is allowed by all to have been a fable.

We think ourselves authorised to say that this attraction of the loadstone for iron, or this tendency of iron to the loadstone, is a secondary phenomenon, and is the consequence of the proper disposition of the induced magnetism. The proofs already given of the compound nature of this phenomenon, namely, that it arises from the excess of two attractions above two repulsions, need (we imagine) no addition. But the following considerations place the matter beyond doubt.

1. The magnetism of the two poles is evidently of an opposite nature; the one repelling what the other attracts. If the one attracts iron, therefore, the other should repel it. But each pole, by inducing a magnetism opposite to its own, on the nearest end of the iron, and the same with its own on the remote end, and its action diminishing with an increase of distance, there must always be an excess of attraction, and the iron must be attracted.

2. Each of the magnets A and B, in either of the positions represented in fig. 12. would alone attract the piece of common iron C. But when placed together, the fourth pole of A tends to render the upper end of C a north pole; while the north pole of B tends to make it a south pole. If their actions be nearly equal, the weight of C cannot be supported by the magnetism induced by any difference of action that may remain. While C is hanging by B alone, let A be gradually brought near; it gradually destroys the action of the north pole of B, so that C gradually loses its magnetism and polarity, and its weight prevails.

§. In all three cases where the induction of magnetism is slow, the attraction is weak in proportion. This is particularly remarked by Dr Gilbert. If we take pieces of common iron, and of steel of different tempers, but all of the same size and form, we shall find that the iron is much more strongly attracted than any of the steel, and that the attraction for each of them is weaker in proportion as they are harder. This diversity is so accurately observed, that when the piece is thoroughly

\subsection*{SUPL (E0070)}
SUPL. Vol. II. Part I.


\section*{Source record}
\begin{description}
\item[Scan ID] \texttt{britannica:309f765b1ecc83e0a7f7762e44b4a927a7555443:train:756}
\item[Image file] \texttt{page-0756-33eb91fcbe8d.jpg}
\item[Image SHA-256] \texttt{33eb91fcbe8d9d080f38ca26a23b791f56382bebff38a428cd4ac3cd30fb48f3}
\item[Dataset] NationalLibraryOfScotland/encyclopaedia\_britannica\_illustrated at revision \texttt{309f765b1ecc83e0a7f7762e44b4a927a7555443} (CC0-1.0)
\item[Extraction] model \texttt{nvidia/nemotron-3.5-parse}, raw reply \texttt{extract/raw/page-0756-33eb91fcbe8d.json}
\item[Raw reply SHA-256] \texttt{88fc3e313b0a120cf1f77e70b9045f86bbca4f346ea6a214529a33efabbc5d12}
\item[Correction] model \texttt{nvidia/Nemotron-3.5-Super-VL}, block crops, raw replies \texttt{correct/raw/crop\_nothink/page-0756-33eb91fcbe8d-b*.json}; parse: 1, super\_vl: 15
\item[Running header] MAGNETISM.
\end{description}
\end{document}
